\section{系统部署与调试建议}

将这套 Compound AI 整体部署到真实产品中，需要在 search、APAC、DSPy 之间建立清晰的接口，并且对每个中间结果添加监控：
\begin{itemize}
    \item Search trace output 需记录 cost/constraint flags，便于 debug；
    \item APAC 的 multi-step question log 需保存用户 intent + agent ask list；
    \item DSPy 的 latent cost \(C\) 和 solver output 应同时上报，以便对齐不同 benchmark。
\end{itemize}

\begin{importantbox}{调试与部署要点}
1）加 monitor：在 inference 中对每次 trace、question、plan 设置 sanity check；2）加缓存：search trace 常常重复，可缓存并复用；3）加 fallback：当 agent judge 指出约束违规时，回退到 hybrid solver 层重新规划；4）加版本控制：每次 DPO 训练都要附带 judge 模型版本。
\end{importantbox}

\begin{knowledgebox}{如何保持链路透明}
在 pipeline 的每个阶段写入 structured log（trace id、agent response、solver plan），并为每个 log 提供 \texttt{trace id -> user request} 的映射；这样即便后台出现崩溃，也能快速回放导致错误的原因。
\end{knowledgebox}

\subsection{本章小结}

部署 Compound AI 时，必须让每个模块的输出可观测、可回放，并在 judge/solver/trace 之间建立版本联动，才能快速定位推理失效并执行回滚。

